\documentclass{article}
\usepackage{amsmath,amssymb}
\usepackage{xurl}
\usepackage[hidelinks]{hyperref}
\setlength{\emergencystretch}{2em}
% Shared commands for notes in docs/paper-gaps/.

\DeclareUnicodeCharacter{2080}{\ifmmode{}_{0}\else\textsubscript{0}\fi}
\DeclareUnicodeCharacter{2081}{\ifmmode{}_{1}\else\textsubscript{1}\fi}
\DeclareUnicodeCharacter{2082}{\ifmmode{}_{2}\else\textsubscript{2}\fi}
\DeclareUnicodeCharacter{2099}{\ifmmode{}_{n}\else\textsubscript{n}\fi}

\newcommand{\Lean}{\textsc{Lean}}
\ExplSyntaxOn
\NewDocumentCommand{\leanid}{m}{
  \ifmmode
    \textsf{\detokenize{#1}}
  \else
    \group_begin:
    \urlstyle{sf}
    \tl_set:Nn \l_tmpa_tl {#1}
    \tl_replace_all:Nnn \l_tmpa_tl {\_} {_}
    \exp_args:NV \path \l_tmpa_tl
    \group_end:
  \fi
}
\ExplSyntaxOff
\providecommand{\tr}{\operatorname{tr}}
\newcommand{\tnleanIssueBase}{https://github.com/LionSR/TNLean/issues/}
\newcommand{\tnleanPrBase}{https://github.com/LionSR/TNLean/pull/}
\newcommand{\tnleanDocBase}{https://sirui-lu.com/TNLean/docs/}
\newcommand{\ghissue}[1]{\href{\tnleanIssueBase#1}{GitHub issue \##1}}
\newcommand{\ghpr}[1]{\href{\tnleanPrBase#1}{PR \##1}}
\newcommand{\doclink}[1]{\href{\tnleanDocBase#1.html}{\path{#1}}}

% Dirac notation and the identity operator, as in blueprint/src/macros/common.tex.
% \providecommand keeps note-local definitions authoritative.
\usepackage{dsfont}
\providecommand{\ket}[1]{|#1\rangle}
\providecommand{\bra}[1]{\langle #1|}
\providecommand{\braket}[2]{\langle #1 | #2 \rangle}
\providecommand{\ketbra}[2]{|#1\rangle\!\langle #2|}
\providecommand{\Id}{\mathds{1}}
\providecommand{\one}{\mathds{1}}
\providecommand{\C}{\mathbb{C}}
\providecommand{\MN}[1]{M_{#1}(\C)}
\providecommand{\MD}{M_D(\C)}

% External referencing for paper sources.  The build helper writes sanitized
% auxiliary files under build/paper-aux/ and excludes duplicated source labels.
\usepackage{xr}

% Import a generated paper auxiliary file with a caller-supplied label prefix.
% The candidate paths support builds from docs/paper-gaps/, the repository root,
% and build/paper-aux/ itself.
\newcommand{\importpaperaux}[2]{%
  \IfFileExists{../../build/paper-aux/#2.aux}{%
    \externaldocument[#1]{../../build/paper-aux/#2}%
  }{%
    \IfFileExists{build/paper-aux/#2.aux}{%
      \externaldocument[#1]{build/paper-aux/#2}%
    }{%
      \IfFileExists{#2.aux}{%
        \externaldocument[#1]{#2}%
      }{}%
    }%
  }%
}

% General external reference macros
\newcommand{\paperref}[2]{\ref{#1-#2}}
\newcommand{\papereqref}[2]{(\ref{#1-#2})}

% CPSV16 source (1606.00608 / MPDO-22-12-17-2.tex) external document and shortcuts
\importpaperaux{cpsv16-}{1606.00608/MPDO-22-12-17-2-tnlean}
\newcommand{\cpsvref}[1]{\paperref{cpsv16}{#1}}
\newcommand{\cpsveqref}[1]{\papereqref{cpsv16}{#1}}

% Verdict marker: \gapnote{<kind>}{<status>}. Kind is the policy
% classification (clarification, local-correction, scope-restriction,
% unfaithful, false-source, open-gap); status is open, wip (elimination
% underway), resolved, or historical. Machine-read by the site index;
% prints a verdict line.
\newcommand{\gapnote}[2]{\par\noindent\textbf{Verdict:} #1 (#2).\par}

\title{The Fibonacci Module: Classification on at Most Two Blocks}
\author{}
\date{2026-09-27}
\begin{document}
\maketitle
\gapnote{scope-restriction}{open}

\section{The source claim}
For a matrix product operator representing the Fibonacci fusion category
$\{1,\tau;\ \tau\times\tau=1+\tau\}$, Garre-Rubio, Lootens and Moln\'ar state
that ``the only solution for an invariant MPS subspace is characterized by two
blocks $\{x_1,x_\tau\}$'' with $\tau\cdot x_1=x_\tau$ and
$\tau\cdot x_\tau=x_1+x_\tau$~\cite{GarreRubio2023MPOSym}.\footnote{Source
    traceability: \path{Papers/2203.12563/REsubmission.tex}, lines 1991--1993.}
The multiplicities $M_{a,x}^y$ of the blocks form a nonnegative integer
representation of the fusion ring,
\begin{align}
    \sum_c N_{ab}^c M_{c,x}^y & =\sum_z M_{b,x}^z M_{a,z}^y ,
    \label{eq:glm23_fib_nimrep}
\end{align}
and the unit acts on the blocks as the identity. For the Fibonacci ring,
(\ref{eq:glm23_fib_nimrep}) at $a=b=\tau$ reads $M_\tau^2=\mathbb 1+M_\tau$ for
the matrix $M_\tau$ with entries $M_{\tau,x}^y$. The claim is stated for
invariant subspaces with any number of blocks.

\section{What is formalized}
The nonnegative integer representations with the unit acting as the identity
are classified on one and on two blocks. On one block, $m^2=1+m$ has no
solution in $\mathbb N$. On two blocks, the entrywise equations of
$M_\tau^2=\mathbb 1+M_\tau$ force the diagonal entries to be $0$ or $1$ with
sum $1$ and both off-diagonal entries to be $1$, so $M_\tau$ is
$\begin{pmatrix}0&1\\1&1\end{pmatrix}$ or its relabelling
$\begin{pmatrix}1&1\\1&0\end{pmatrix}$: the regular representation. Combined
with integrality of the multiplicities of normal blocks, a symmetric family of
at most two normal blocks with the unit acting trivially has exactly two
blocks, transforming as in the source.\footnote{Results:
    \leanid{FibonacciCompression.not_isNIMRep_fibNim_of_card_eq_one},
    \leanid{FibonacciCompression.exists_equiv_of_isNIMRep_fibNim_of_card_eq_two}
    and
    \leanid{FibonacciCompression.exists_equiv_of_isMPOSymmetricFamily_fibBlock}.}

\section{Boundary}
Not formalized: families with three or more blocks. For these the source
claim amounts to the statement that every indecomposable nonnegative integer
representation of the Fibonacci ring with the unit acting trivially is the
regular one, which rests on the classification of module categories over the
Fibonacci category that the source cites. A decomposable family, a direct
sum of copies of the regular representation, also satisfies
(\ref{eq:glm23_fib_nimrep}); the source's uniqueness is up to this
decomposition. The periodic-boundary reading of the symmetry hypothesis is
recorded separately in the note on the periodic-boundary layer.

Elimination plan: show that a nonnegative integer matrix $M$ with
$M^2=\mathbb 1+M$ is, after a permutation of the blocks, a direct sum of
copies of $\begin{pmatrix}0&1\\1&1\end{pmatrix}$. Order the strongly connected
classes of the support of $M$ so that $M$ is block upper triangular. Each
diagonal block $F$ is irreducible with $F^2=\mathbb 1+F$, so its eigenvalues
are $\varphi$, simple by the Perron--Frobenius theorem, and $-\varphi^{-1}$;
the trace $\varphi-(k-1)\varphi^{-1}$ of a block of size $k$ is an integer
only for $k=2$. An off-diagonal block $X$ between classes with diagonal
blocks $F$ and $F'$ satisfies $FX+XF'-X\leq0$ entrywise, and pairing with
positive Perron--Frobenius eigenvectors $uF=\varphi u$, $F'w=\varphi w$ gives
$\sqrt5\,uXw\leq0$, hence $X=0$.
\bibliographystyle{alpha}
\bibliography{references}
\end{document}
